\par
\begin{table}[h]\centering\small
\caption{Complete residual classes at $\xs$. Values are $(n_*\cdot p-d_*)/\sqrt3$; every nonzero row is strictly negative.}\label{tab:residuals}
\begin{tabular}{@{}ll@{}}\toprule
Exact residual & Branch labels\\\midrule
$-1 + \frac{\sqrt{2}}{2}$ & \br{E06-}, \br{E08+}, \br{E15-}, \br{E18+} \\
$- \frac{1}{2} + \frac{\sqrt{2}}{4}$ & \br{E07-}, \br{E14+} \\
$- \frac{3}{2} + \frac{\sqrt{2}}{2}$ & \br{E01-}, \br{E02-}, \br{E03-} \\
$1 - \frac{3 \sqrt{2}}{4}$ & \br{E19+}, \br{E19-} \\
$\frac{1}{2} - \frac{\sqrt{2}}{2}$ & \br{E04+}, \br{E11-}, \br{E12+}, \br{E20-} \\
$2 - 2 \sqrt{2}$ & \br{E04-}, \br{E11+}, \br{E12-}, \br{E20+} \\
$3 - \frac{5 \sqrt{2}}{2}$ & \br{E01+}, \br{E02+}, \br{E03+} \\
$3 - \frac{9 \sqrt{2}}{4}$ & \br{E09-}, \br{E17-} \\
$\frac{3}{2} - \frac{3 \sqrt{2}}{2}$ & \br{E05-}, \br{E10+}, \br{E16+}, \br{E21-} \\
$\frac{3}{2} - \frac{5 \sqrt{2}}{4}$ & \br{E13+}, \br{E13-} \\
$\frac{5}{2} - 2 \sqrt{2}$ & \br{E07+}, \br{E14-} \\
$\frac{7}{2} - \frac{5 \sqrt{2}}{2}$ & \br{E06+}, \br{E08-}, \br{E15+}, \br{E18-} \\
$0$ & \br{E05+}, \br{E09+}, \br{E10-}, \br{E16-}, \br{E17+}, \br{E21+} \\
\bottomrule\end{tabular}
\src{Accepted candidate residual table~\cite{investigation}. For example, $(7-5\sqrt2)/2<0$ because $49<50$.}
\end{table}\par
